% Part: sets-functions-relations
% Chapter: infinite
% Section: dedekind-induction

\documentclass[../../../include/open-logic-section]{subfiles}

\begin{document}

\olfileid{sfr}{infinite}{induction}
\olsection[Volledige inductie]{Dedekindalgebra's en volledige inductie}

Nu komt het belangrijkste punt: een Dedekindalgebra---\emph{elke}
Dedekindalgebra---kan de rol van de natuurlijke getallen overnemen. Dat
volgt meteen uit dit eenvoudige resultaat:

\begin{thm}[Volledige inductie]\ollabel{thm:dedinfiniteinduction}
Laat $N, s, o$ samen een Dedekindalgebra vormen. Neem een willekeurige
verzameling $X$. Dan geldt:
	\begin{center}
		als $o \in X$ en $(\forall {n} \in N \cap X){s}({n}) \in X$, {dan} $N \subseteq X$.
	\end{center}
\end{thm}

\begin{proof}
Volgens de definitie van een Dedekindalgebra is $N =
\closureofunder{s}{o}$. Stel nu dat ${o} \in X$ en $(\forall {n} \in
N)(n \in X \lif {s}({n}) \in X)$. Het gekozen element ligt ook in
de verzameling van de algebra, en dus geldt ${o} \in N \cap X$.
Neem nu ${n} \in N \cap X$.
Volgens de aanname geldt ${s}({n}) \in X$. Ook geldt
${s}({n}) \in N$, omdat $s \colon N \to N$. Dus is $N \cap X$
$s$-gesloten. Deze verzameling bevat het gekozen element. Zij bevat dus
ook de kleinste gesloten verzameling met dat element. Daarom is $N =
\closureofunder{s}{o} \subseteq X$.
\end{proof}

Inductie is een kenmerkende eigenschap van de natuurlijke getallen.
Daarom kunnen we van elke Dedekind-oneindige verzameling een
Dedekindalgebra maken en die algebra de rol van de natuurlijke getallen
laten spelen.

\olref{thm:dedinfiniteinduction} beschrijft inductie met
\emph{verzamelingstheoretische} begrippen. We kunnen dezelfde regel
ook in een bekende vorm opschrijven:

\begin{cor}\ollabel{natinductionschema}
Laat $N, s, o$ samen een Dedekindalgebra vormen. Neem een formule
$\phi(x)$ waarin parameters mogen voorkomen. Dan geldt:
\begin{center}
	als $\phi(o)$ en $(\forall {n} \in N)(\phi(n)\lif
	\phi({s}({n})))$, {dan} $(\forall n \in N)\phi(n)$
\end{center}
\end{cor}

\begin{proof}
Neem $X = \Setabs{n \in N}{\phi(n)}$ en gebruik nu
\olref{thm:dedinfiniteinduction}.
\end{proof}

Hierboven zeiden we dat een formule ``parameters mag hebben''. Grofweg
betekent dit dat we voor elke keuze van objecten $c_1, \ldots, c_k$
mogen werken met $\phi(x, c_1, \ldots, c_k)$. We kunnen hetzelfde ook
precies opschrijven zonder het woord ``parameters''. Neem een formule
$\phi(x, v_1, \ldots, v_k)$ en schrijf al haar vrije variabelen erbij.
Dan geldt:
	\begin{align*}
			\forall v_1 \ldots \forall v_k((&\phi(o, v_1,\ldots, v_k) \land {}\\
			&	(\forall x \in N)(\phi(x,v_1, \ldots, v_k) \lif \phi(s(x), v_1,\ldots, v_k))) \lif {}\\
			&\hspace{3em} (\forall x \in N)\phi(x, v_1,\ldots, v_k))
	\end{align*}
De woorden ``met parameters'' maken dit duidelijk veel makkelijker te
lezen. (In \olref[sth][][]{part} zullen we deze manier van spreken vaak
gebruiken.)

Terug naar Dedekindalgebra's. Voor elke Dedekindalgebra kunnen we ook
de gewone rekenfuncties voor optellen, vermenigvuldigen en
machtsverheffen vastleggen. Dat is niet eenvoudig. We hebben daarvoor
de techniek van de \emph{recursieve definitie} nodig. Die techniek
leggen we pas veel later uit en rechtvaardigen we daar in een veel
algemenere situatie. (Wie nieuwsgierig is, kan dit deel na
\olref[sth][ord-arithmetic][]{chap} nog eens lezen, of een andere uitleg
bekijken, zoals \citealt[pp.~95--8]{Potter2004}.) Als $N, s, o$ samen
een Dedekindalgebra vormen, kunnen we uiteindelijk deze regels
vastleggen:
\begin{align*}
	{a} + {o} &= {a} & & & {a} \times {o} &= {o} & & & {a}^{o} &= s(o)\\
{a} + {s}({b}) &= {s}({a}+{b}) &&& {a} \times {s}({b}) &= ({a}\times {b}) + {a}  & & & {a}^{{s}({b})} &= {a}^{b} \times {a}
\end{align*}
Daarna kunnen we laten zien dat deze rekenfuncties werken zoals we
verwachten.

\end{document}
